\chapter{\babSatu}
\label{bab:pendahuluan}
\hyphenation{FastAPI SQLite statsmodels mlxtend ihelon Starbucks Laravel ba-ris-ta ta-kar-an pe-nye-duh-an pen-ska-la-an re-ko-men-da-si in-ven-ta-ris per-se-di-a-an pe-ra-mal-an stan-dar-di-sa-si o-pe-ra-si-o-nal kon-sis-ten-si eks-trak-si pe-la-tih-an es-pre-so pa-sang-an peng-a-da-an pe-ngem-bang-an ke-ter-kait-an pem-ban-ding di-ka-li-kan pe-me-rik-sa-an ke-dai dataset back-end front-end frame-work end-point re-stock}

% Sumber fakta bab ini: tools/PANDUAN_PENULISAN.md, app/SPEC.md, "Skripsi Arahan Pak Budi.pdf",
% referensi_terverifikasi.json (klaim dicocokkan dengan abstract_evidence), lembar data dataset di
% dataset/raw/*/DATASHEET.md, dan results/external/E5_* (hanya pola shot menu Starbucks yang bersifat
% deskriptif). Bab ini tidak memuat angka kinerja hasil evaluasi (hasil ada di Bab 4).

\section{Latar Belakang}
\label{sec:latar-belakang}

Konsumsi kopi domestik Indonesia diperkirakan masih meningkat. Laporan \textit{Coffee Annual} untuk Indonesia yang disusun staf Departemen Pertanian Amerika Serikat memperkirakan konsumsi kopi domestik pada periode 2026/2027 sebesar 4,83 juta karung berukuran 60~kg, naik 20.000 karung dari periode sebelumnya \citep*{rahmanulloh2026coffee}. Laporan yang sama mencatat bahwa gerai kopi di ruang publik terus melayani konsumen kelas atas, termasuk pelanggan generasi Z, dan bahwa kopi susu tetap menjadi minuman pengantar yang populer dan tersedia luas di berbagai kedai kopi. Pertumbuhan industri kedai kopi (\textit{coffee shop}) Indonesia yang pesat juga menjadi titik tolak penelitian peramalan penjualan kedai kopi \citep*{azzahra2025forecasting}. \citet*{pattisealnno2026comparative} mengaitkan naiknya konsumsi kopi di Indonesia dengan risiko kelebihan dan kekurangan stok yang dihadapi kedai kopi. Dengan demikian, kedai kopi dituntut menyajikan minuman yang baik sekaligus mengelola bahan bakunya dengan cermat.

Arahan penelitian ini menempatkan konsistensi rasa dan kecepatan pelayanan sebagai kunci utama kepuasan pelanggan kedai kopi, dan penelitian di Indonesia menunjukkan pentingnya mutu produk dan layanan bagi pelanggan kedai kopi. Pada sebuah kedai kopi di Bali, kualitas produk dan kualitas layanan berpengaruh positif signifikan terhadap kepuasan dan niat beli ulang pelanggan, dengan kepuasan sebagai mediator \citep*{natalia2023role}. \citet*{candra2022impact} menemukan bahwa pengalaman sensorik pengunjung kedai kopi di Indonesia memengaruhi emosi positif dan negatif serta niat perilaku mereka, lalu menyarankan pemilik merancang setiap aspek suasana sensorik agar pengunjung kembali. Pada tingkat prosedur, unsur prosedur operasional standar (\textit{standard operating procedure}, SOP) yang efisien, efektif, dan konsisten masing-masing berpengaruh signifikan terhadap kualitas pelayanan sebuah kafe di Jember menurut persepsi pelanggannya \citep*{chapitoline2024pengaruh}. Temuan-temuan ini mengisyaratkan bahwa rasa yang berubah-ubah antarkunjungan dan pelayanan yang lambat berisiko melemahkan kepuasan dan niat kembali pelanggan.

Menjaga rasa tetap sama dari cangkir ke cangkir tidak mudah karena mutu seduhan dipengaruhi banyak variabel. Tinjauan \citet*{cordoba2020coffee} menunjukkan bahwa metode seduh berbeda dalam tekanan ekstraksi, rasio kopi dan air, kualitas air, waktu kontak, distribusi ukuran partikel, dan suhu, dan seluruh variabel tersebut mengubah ekstraksi senyawa yang membentuk cita rasa. Pada kopi seduh tetes, kekuatan seduhan (\textit{total dissolved solids}, TDS) dan persentase ekstraksi mengubah profil sensori secara sistematis: 21 dari 30 atribut sensori berbeda nyata antarseduhan, dan rasa manis berbanding terbalik dengan TDS \citep*{frost2020effects}. Temuan semacam ini dirangkum \citet{guinard2023new} menjadi versi baru \textit{Coffee Brewing Control Chart} yang mengaitkan atribut sensori dan kesukaan konsumen dengan kekuatan seduhan, rendemen ekstraksi, dan rasio seduh. Espreso bahkan disebut sebagai format kopi yang paling rentan terhadap variasi, dan penurunan rendemen pada gilingan halus diduga kuat berasal dari aliran yang tidak merata, yang menurunkan reprodusibilitas serta membuang bahan baku \citep{cameron2020systematically}. Faktor manusia juga berperan: metode seduh filter memakan waktu dan bergantung pada keterampilan barista, dan tekanan serta turbulensi saat barista menuang air menjadi variabel kunci reprodusibilitas ekstraksi, walaupun seduhan V60 dan Pure Brew dari enam barista ahli dalam studi yang sama menunjukkan reprodusibilitas TDS dan waktu seduh yang sangat baik \citep{santanatoglia2023impact}. Penulis menafsirkan temuan terakhir ini sebagai petunjuk bahwa konsistensi setingkat barista ahli dapat dijadikan sasaran prosedur seduh yang terstandar bagi barista yang belum berpengalaman. \citet{park2023robot} menemukan bahwa kopi seduhan barista manusia memiliki variasi senyawa volatil yang lebih besar daripada seduhan robot barista, walaupun penerimaan konsumen terhadap keduanya tidak berbeda nyata.

Arahan penelitian ini mencatat tiga kendala yang sering muncul dalam operasional harian kedai kopi. Pertama, takaran resep tidak konsisten akibat pergantian \textit{shift} barista. Kedua, pelatihan staf baru memerlukan waktu yang lama. Ketiga, manajemen inventaris bahan baku seperti biji kopi, susu, dan sirup belum terintegrasi dengan baik. Arahan yang sama meminta logika perhitungan dinamis untuk menyesuaikan takaran resep terhadap variasi pesanan. Tanpa dukungan sistem, penyesuaian itu harus dihitung ulang secara manual oleh barista. Kendala serupa tercatat pada studi kasus manajemen risiko berbasis ISO 31000 di Kanca Coffee, Bogor, yang menempatkan malfungsi mesin espreso, pergantian karyawan (\textit{turnover}), dan gangguan pasokan bahan baku sebagai risiko operasional paling kritis, lalu mengusulkan mitigasi antara lain program pelatihan dan retensi karyawan serta perbaikan prosedur manajemen persediaan \citep*{akbar2026operational}.

\textit{Turnover} karyawan yang tinggi merupakan realitas yang dikenal luas di industri \textit{hospitality} \citep*{dogru2023employee}. \citet*{lee2022barista} menggambarkan pekerjaan barista sebagai jalinan kerumitan operasional yang melampaui sekadar membuat kopi, dengan keterampilan yang berkembang sebagai kebiasaan (\textit{habitual powers}). Penulis menafsirkan kedua temuan tersebut sebagai petunjuk bahwa pengetahuan takaran yang hanya tersimpan dalam kebiasaan barista berpengalaman mudah hilang ketika barista berganti \textit{shift} atau keluar, sehingga pengetahuan itu perlu dikodifikasi. \citet*{naumov2023employment} berpendapat bahwa resep standar memungkinkan pelatihan karyawan jasa makanan yang lebih cepat, walaupun argumen tersebut bersifat konseptual. Literatur pelatihan industri menyebut alat bantu kerja (\textit{job performance aids}), seperti panduan prosedural dan daftar periksa, sebagai alternatif atau pelengkap pelatihan yang hemat biaya selama alat tersebut mengurangi beban ingatan dan merinci tindakan beserta waktunya \citep*{campbell1996job}. Dalam domain kopi, sistem realitas virtual cerdas yang memberi umpan balik korektif meningkatkan kinerja belajar 103 mahasiswa yang berlatih menyeduh \textit{pour over} \citep*{yu2021developing}, tetapi karena subjeknya mahasiswa dan perangkatnya realitas virtual, penelitian ini tidak mengklaim bahwa aplikasinya memperpendek masa pelatihan barista tanpa bukti dari uji penerimaan penggunanya sendiri.

Kendala inventaris juga banyak dilaporkan pada kedai kopi di Indonesia. Kekurangan bahan baku mengganggu operasional dan menurunkan kepuasan pelanggan, sedangkan pembelian mendadak membuat operasional tidak stabil dan menaikkan biaya persediaan \citep*{nurhalizapratiwi2025analisis}. Di Kava Coffee Shop, Kudus, pencatatan di atas kertas menyebabkan data tidak akurat, pemesanan ulang terlambat, serta kehabisan dan kelebihan stok \citep*{taufiq2025implementation}. Pencatatan manual dengan buku besar pada sebuah kafe di Pontianak menimbulkan kesalahan pencatatan, pemeriksaan stok yang lambat, dan risiko dokumen hilang \citep*{abdur2026sistem}. Sebaliknya, penumpukan bahan baku kopi pada sebuah kedai kopi menimbulkan kerugian finansial karena bahan rusak selama penyimpanan \citep*{nasution2025application}. Selama pemakaian bahan tidak terhubung dengan resep dan produksi, stok di sistem hanya dapat diperbarui melalui pencatatan terpisah yang rawan terlambat dan keliru.

Resep standar adalah jawaban klasik atas masalah konsistensi. \citet*{naumov2023employment} mendefinisikan resep standar sebagai panduan untuk menyiapkan makanan atau minuman secara konstan pada mutu yang diharapkan, dengan unsur nama resep, hasil (\textit{yield}), ukuran porsi, berat atau jumlah, satuan takar, bahan, instruksi, dan catatan, sehingga mutu, kuantitas, dan biaya porsi dapat diprediksi dan dikendalikan. Agar resep juga dapat mengendalikan persediaan, komposisinya perlu diperlakukan sebagai daftar bahan (\textit{bill of materials}, BOM). \citet{kamal2021bill} mengusulkan analisis BOM menu karena analisis manual atas bahan penyusun menu yang kompleks tidak efisien. Dengan BOM, setiap cangkir yang dicatat dapat diterjemahkan langsung menjadi pemakaian biji kopi, susu, dan sirup pada buku besar stok. Bentuk digital juga membantu: pada tugas perbaikan bilah turbin, instruksi kerja digital berbantuan \textit{augmented reality} dipilih oleh kesepuluh pekerja yang diuji dan membuat tugas selesai 21\% lebih cepat dengan beban kerja yang dirasakan 26\% lebih rendah daripada dengan dokumen kertas \citep*{eversberg2023evaluating}, walaupun konteksnya berbeda dari bar kopi.

Resep standar ditulis untuk satu ukuran sajian, padahal pesanan di bar datang dalam berbagai variasi. Cara yang paling umum untuk menyesuaikan resep adalah metode faktor konversi, yaitu membagi hasil yang dibutuhkan dengan hasil resep, lalu mengalikan setiap bahan dengan faktor tersebut \citep*{bccook2015basic}. Buku ajar yang sama mengingatkan bahwa bumbu harus dinaikkan dengan hati-hati karena melipatduakan atau melipattigakannya dapat berdampak buruk, waktu pengolahan dapat berubah dengan peralatan yang berbeda, dan penyesuaian halus hanya dapat dipelajari dari pengalaman karena tidak ada aturan baku, sehingga resep hasil konversi perlu diuji. Buku ajar kuliner \citet*{gisslen2025professional} bahkan menyediakan bagian khusus tentang masalah dalam konversi resep.

Di bar kopi, penskalaan linier menghadapi tiga batasan praktis. Pertama, espreso diekstrak dalam satuan \textit{shot} utuh, sehingga faktor pengali 1,5 tidak dapat dijalankan sebagai satu setengah \textit{shot}; takaran per \textit{shot} juga tidak dapat diubah sembarangan karena perbedaan massa total minuman espreso yang diekstrak berpengaruh jauh lebih besar terhadap massa komponen di dalam cangkir daripada perubahan gilingan, laju alir, dan suhu \citep*{schmieder2023influence}. Kedua, volume dibatasi kapasitas cangkir: arahan menetapkan cangkir 200--250~ml untuk Caffè Latte dan 150--180~ml untuk Cappuccino, sedangkan Ice Brown Sugar Shaken Espresso diisi susu hingga penuh (\textit{top up}), sehingga tambahan \textit{shot} semestinya mengurangi susu alih-alih membuat minuman meluap. Ketiga, seduhan manual bergantung pada rasio seduh (\textit{brew ratio}). Resep V60 pada arahan memakai 15~g kopi untuk 225~ml air (1:15) dengan target tuang kumulatif 45, 130, dan 225~ml, sehingga dosis 20~g menuntut target 60; 173,3; dan 300~ml agar rasionya tetap. Pada seduhan imersi penuh, kekuatan seduhan kira-kira berbanding terbalik dengan rasio air terhadap kopi sehingga rasio seduh menjadi parameter terpenting \citep*{liang2021equilibrium}; bagi V60 yang merupakan seduhan perkolasi, hubungan ini hanya berlaku sebagai pendekatan.

% sumber: results/external/E5_scaling_starbucks.json (analysis_sets.A_integer_shots: 7 minuman;
% results.A_integer_shots.nominal_floz.shot_patterns: 1-1-2-2 = 6 minuman, 1-2-3-4 = 1 minuman) dan
% results/external/E5_implied_shots.csv (kafein dan shot tersirat Caffe Latte per ukuran).
% Dataset: tatman2017nutrition; volume nominal cangkir (bukan kolom dataset): dimberg2023starbucks.
Pola ukuran pada menu sebuah jaringan kedai kopi besar memperlihatkan bahwa praktik penyajian memang tidak linier. \textit{Dataset} publik Nutrition Facts for Starbucks Menu di Kaggle \citep*{tatman2017nutrition} mencantumkan kandungan kafein minuman untuk setiap ukuran, dan dengan kafein satu \textit{shot} espreso sebesar 75~mg pada \textit{dataset} yang sama, jumlah \textit{shot} setiap ukuran dapat diturunkan. Tabel~\ref{tab:pola-shot} menunjukkan bahwa Caffè Latte ukuran Short dan Tall sama-sama memakai satu \textit{shot}, sedangkan ukuran Grande dan Venti memakai dua \textit{shot}. Penskalaan proporsional dari ukuran Grande justru menghasilkan 1,5 \textit{shot} untuk Tall dan 2,5 \textit{shot} untuk Venti. Dari tujuh minuman berbasis espreso yang dijual dalam empat ukuran panas dan kafeinnya merupakan kelipatan utuh 75~mg, enam mengikuti pola bertingkat seperti ini (1, 1, 2, dan 2 \textit{shot} dari Short hingga Venti), sedangkan Caffè Americano bertambah satu \textit{shot} pada setiap ukuran (1, 2, 3, dan 4 \textit{shot}). Pola ini hanya menggambarkan menu satu jaringan kedai di Amerika Serikat sekitar tahun 2017 dan dapat berbeda dengan kedai kopi di Indonesia, tetapi cukup untuk menunjukkan bahwa algoritma penskalaan perlu memperlakukan \textit{shot} espreso sebagai bahan diskret, sedangkan susu, air, dan sirup dapat tetap diperlakukan sebagai bahan kontinu.

% sumber: results/external/E5_implied_shots.csv (kafein, shot tersirat); shot proporsional = 2 x volume / 16 fl oz.
\begin{table}[htbp]
\centering
\begin{threeparttable}
\caption{Kafein dan Jumlah \textit{Shot} Caffè Latte per Ukuran pada \textit{Dataset} Menu Starbucks}
\label{tab:pola-shot}
\begin{tabular}{lcccc}
\toprule
Ukuran & \begin{tabular}[c]{@{}c@{}}Volume Nominal\\(fl oz)\tnote{a}\end{tabular} & \begin{tabular}[c]{@{}c@{}}Kafein\\(mg)\end{tabular} & \begin{tabular}[c]{@{}c@{}}\textit{Shot}\\Tersirat\tnote{b}\end{tabular} & \begin{tabular}[c]{@{}c@{}}\textit{Shot} Proporsional\\dari Grande\end{tabular} \\
\midrule
Short  & 8  & 75  & 1 & 1,00 \\
Tall   & 12 & 75  & 1 & 1,50 \\
Grande & 16 & 150 & 2 & 2,00 \\
Venti  & 20 & 150 & 2 & 2,50 \\
\bottomrule
\end{tabular}
\begin{tablenotes}\footnotesize
\item[a] Volume nominal minuman panas dari sumber sekunder, yaitu \citet*{dimberg2023starbucks}, bukan dari \textit{dataset}.
\item[b] Kafein dibagi 75~mg, yaitu kafein Espresso Solo (satu \textit{shot}) pada \textit{dataset} yang sama.
\end{tablenotes}
\end{threeparttable}
\end{table}

Penelitian komputasi resep yang paling dekat berasal dari luar bar kopi. \citet*{benita2026web} membangun platform web yang menskalakan resep masakan rumah tangga per porsi, menyediakan penggantian bahan dan pengaturan tingkat pedas, serta asisten memasak berbasis kecerdasan buatan; abstraknya menyajikan prototipe dan rencana evaluasi tanpa hasil evaluasi. \citet{choi2023kitchenscale} memprediksi kuantitas dan satuan bahan dari konteks resep dengan model bahasa praterlatih, sedangkan \citet*{moralesgarzon2025service} mengadaptasi resep melalui penggantian bahan sesuai preferensi, tujuan kesehatan dan keberlanjutan, serta alergi pengguna. Berdasarkan abstraknya, ketiga penelitian tersebut tidak membahas batasan khas bar kopi seperti \textit{shot} utuh, kapasitas cangkir, dan pengisian susu hingga penuh, serta tidak membandingkan hasilnya dengan penskalaan linier. Penelitian ini mengisi kekosongan tersebut dengan algoritma penskalaan berbasis batasan yang deterministik dan dapat diaudit, lalu membandingkannya dengan penskalaan linier sebagai pembanding.

Selain resep dan stok, kedai kopi yang mencatat transaksi secara digital mengumpulkan data penjualan yang dapat dimanfaatkan kecerdasan buatan (\textit{artificial intelligence}, AI). Sistem rekomendasi didefinisikan sebagai perangkat lunak dan teknik yang memberi saran tentang \textit{item} yang paling mungkin menarik bagi pengguna tertentu \citep*{ricci2021recommender}. Dukungan semacam ini relevan karena usaha makanan dan minuman umumnya dikelola langsung oleh pemiliknya dengan akses terbatas pada teknologi peramalan modern, dan manajernya sering meramal permintaan dengan heuristik berbasis pengalaman \citep*{groene2024introduction}. Penelitian ini memanfaatkan peluang tersebut pada tiga keputusan operasional, yaitu pasangan menu yang ditawarkan, jumlah bahan baku yang perlu dipesan ulang, dan koreksi parameter seduh ketika hasil seduhan menyimpang.

Keputusan pertama dapat didukung aturan asosiasi. \citet*{agrawal1993mining} memperkenalkan penambangan aturan asosiasi antar-\textit{item} dari basis data transaksi pelanggan, dan algoritma FP-Growth menambang seluruh pola sering tanpa membangkitkan kandidat sehingga dilaporkan sekitar satu orde besaran lebih cepat daripada Apriori \citep*{han2000mining}. Di Indonesia, FP-Growth sudah diterapkan pada data penjualan Cascara Coffee untuk strategi \textit{cross-selling} dan \textit{up-selling} \citep*{fadillah2021data}, Cafe Over Limit \citep*{hasan2023utilization}, dan Kedai Kopi Raja Palembang \citep*{pratama2025food}, serta pada \textit{dataset} publik The Bread Basket \citep*{mittal2020bread} untuk rekomendasi menu pada sistem kasir \citep*{sihombing2026analisis}. Penelitian-penelitian tersebut menilai aturan dengan \textit{support}, \textit{confidence}, dan \textit{lift}, tetapi sejauh yang dapat ditelusuri dari abstraknya tidak menguji ketepatan rekomendasi pada transaksi periode berikutnya dan tidak membandingkannya dengan metode pembanding sederhana. Padahal, evaluasi yang mengabaikan urutan waktu membuat model belajar dari interaksi yang belum tersedia saat prediksi, sehingga akurasi dan urutan peringkat model menjadi tidak dapat diandalkan \citep*{ji2023critical}. Pembanding sederhana juga perlu disertakan karena algoritma nonpersonal yang naif dapat mengungguli sebagian pendekatan rekomendasi yang umum \citep*{cremonesi2010performance}.

Keputusan kedua bergantung pada peramalan permintaan: peramalan penjualan menu yang tepat memungkinkan pemesanan stok bahan secara presisi dan menekan limbah pangan sebelum konsumsi \citep*{posch2022bayesian}, sedangkan peramalan yang buruk disebut sebagai penyebab utama limbah pangan pada layanan katering \citep*{rodrigues2024machine}. Pada kedai kopi, \citet*{azzahra2025forecasting} memperoleh galat terendah dari ARIMA dibandingkan dengan rata-rata bergerak dan SARIMA untuk penjualan harian satu produk, dan \citet*{pattisealnno2026comparative} menyimpulkan bahwa model statistik sederhana lebih cocok daripada \textit{deep learning} bagi usaha ritel kecil dengan riwayat data terbatas. Namun, ramalan penjualan menu belum langsung menjawab berapa banyak biji kopi atau susu yang harus dipesan. \citet*{goltsos2022inventory} mencatat bahwa penelitian peramalan sering mengabaikan cara ramalan diubah menjadi keputusan pengisian ulang, sedangkan teori persediaan sering menganggap permintaan sudah diketahui. Sistem persediaan kedai kopi yang memakai \textit{safety stock} dan \textit{reorder point} (ROP) pun tidak memakai peramalan permintaan maupun data resep \citep*{taufiq2025implementation}. Penelitian yang paling dekat menggabungkan keduanya dengan cara lain: \citet*{lauren2026optimizing} memasukkan galat ramalan ke perhitungan \textit{safety stock} UMKM kuliner dengan data yang jarang, sedangkan \citet*{piewthongngam2026aienabled} mengubah ramalan LightGBM menjadi kebutuhan bahan baku melalui rasio pemakaian historis terhadap penjualan pada jaringan restoran prasmanan. Penelitian ini menurunkan kebutuhan bahan langsung dari resep standar yang tersimpan di aplikasi.

Keputusan ketiga terjadi di meja bar ketika seduhan terasa terlalu asam, pahit, atau encer. Standar Golden Cup menetapkan kekuatan seduhan 1,15--1,35\% TDS dengan rendemen ekstraksi 18--22\% \citep*{scaa2015golden}, tetapi preferensi konsumen tersebar pada rentang kekuatan dan rendemen yang lebar sehingga titik ``ideal'' pada bagan kendali seduh perlu ditinjau ulang \citep*{cotter2020consumer}. Pendekatan AI yang sudah diajukan untuk parameter seduh antara lain pembelajaran penguatan (\textit{reinforcement learning}) untuk mengoptimalkan variabel \textit{pour over} dari data barista dan penyeduh rumahan di Cianjur \citep*{bramantoro2025optimization} serta \textit{proximal policy optimization} dengan \textit{curriculum learning} \citep*{jeffery2025optimizing}. Sistem pakar dengan \textit{forward chaining} juga telah dipakai di domain kopi untuk memilih biji kopi pascasangrai \citep*{anggraeni2024sistem} dan mendiagnosis kerusakan mesin espreso \citep*{rizal2024sistem}. Penelitian ini memilih sistem pakar berbasis aturan yang bersumber dari bagan kendali seduh, literatur penyeduhan, dan standar resep yang tersimpan di aplikasi, karena rekomendasi berbasis pengetahuan memanfaatkan pengetahuan domain secara eksplisit, berbeda dengan \textit{collaborative filtering} dan \textit{content-based filtering} \citep{uta2024knowledgebased}. Dengan pendekatan ini, setiap saran koreksi dapat ditelusuri ke aturan dan sumber yang memicunya.

Dari sisi sistem informasi, penelitian 2021--2026 di Indonesia, termasuk pada kedai kopi dan usaha kuliner, umumnya berfokus pada satu kelompok kebutuhan. Contohnya adalah sistem stok bahan baku berbasis Laravel \citep*{abdur2026sistem}, sistem persediaan dengan \textit{safety stock} dan ROP \citep*{taufiq2025implementation}, sistem kasir kedai kopi \citep*{rusmara2025pengembangan}, aplikasi kasir berbasis \textit{Progressive Web App} untuk UMKM \citep*{prayudha2024rancang}, sistem harga pokok produksi berbasis komposisi resep \citep*{kartono2025perancangan}, dan digitalisasi SOP pada perusahaan logistik \citep*{shabrina2026sistem}. Sistem yang paling mendekati adalah rekomendasi menu katering industri yang memadukan penyaringan berdasarkan stok dan anggaran, pengelompokan K-Means, dan Apriori, serta memotong stok bahan secara otomatis \citep*{pratama2026automatic}. Sejauh penelusuran penulis, belum ditemukan aplikasi kedai kopi yang sekaligus (1) menyimpan resep standar beserta parameter seduh dan SOP, (2) menghitung ulang takaran dengan batasan praktis bar dan membandingkannya dengan penskalaan linier, (3) memotong stok dari produksi yang dicatat melalui BOM resep, dan (4) menyediakan rekomendasi AI yang dievaluasi terhadap pembanding sederhana dengan pembagian data berbasis waktu, baik pada data sampel fiktif maupun pada data nyata. Celah inilah yang diisi penelitian ini, sedangkan perbandingan rinci dengan penelitian terkait disajikan pada Subbab~\ref{sec:penelitian-terkait}.

Evaluasi modul berbasis transaksi memerlukan data yang tersedia secara terbuka, tetapi dua \textit{dataset} kedai kopi publik yang katalog produknya paling cocok dengan menu aplikasi, yaitu IBM Coffee Shop Sample Data \citep*{chang2019coffee} dan Maven Coffee Shop Sales \citep*{mavenanalytics2023coffee}, merupakan data sampel fiktif menurut penerbitnya masing-masing. Karena itu, kedua \textit{dataset} hanya dipakai sebagai data pengembangan, sedangkan kinerja pasangan menu, peramalan, dan \textit{restock} diperiksa ulang pada data nyata (Subbab~\ref{sec:ruang-lingkup}). Evaluasi ketiga modul itu memakai pembagian data berbasis waktu untuk menghindari kebocoran data, yang ditemukan di 17 bidang ilmu dan pada sebagian kasus menghasilkan kesimpulan yang terlalu optimistis \citep*{kapoor2023leakage}. Dengan cara ini, model hanya dinilai pada data yang terjadi setelah data latihnya, sebagaimana situasi pemakaian yang sebenarnya. Asisten tanya-jawab dievaluasi terpisah pada tolok ukur internal dan tidak diklaim tervalidasi pada barista kedai.

Dari sisi teknologi, arahan penelitian ini mengusulkan \textit{Progressive Web App} (PWA) dengan \textit{framework} PHP Laravel \citep*{laravel2026installation} pada sisi \textit{backend} dan Vue.js pada sisi \textit{frontend}, termasuk \textit{service worker} agar aplikasi tetap dapat dipakai secara luring ketika koneksi internet di kedai tidak stabil; PWA sendiri dipandang sebagai kemungkinan langkah menuju pengembangan lintas platform dari satu basis kode \citep*{majchrzak2018progressive}. Penelitian ini secara sadar menyimpang dari arahan tersebut: aplikasi dibangun dengan FastAPI (Python) dan Vue.js sebagai aplikasi web responsif tanpa kemampuan PWA. Alasannya, modul AI dan pengolahan \textit{dataset} dikerjakan dengan ekosistem Python, yaitu pandas \citep*{mckinney2010data}, scikit-learn \citep{pedregosa2011scikitlearn}, statsmodels \citep*{seabold2010statsmodels}, dan mlxtend \citep*{raschka2018mlxtend}, sehingga API dan AI cukup ditulis dalam satu bahasa, tidak terpisah seperti pada sistem katering \citet*{pratama2026automatic} yang memakai Laravel untuk aplikasinya dan Python untuk komponen pembelajaran mesinnya. FastAPI membangun API dengan Python berdasarkan anotasi tipe standar dan menghasilkan dokumentasi OpenAPI secara otomatis \citep*{ramirez2026fastapi}, dan uji tolok ukur \citet*{bednarz2025benchmarking} menempatkannya di atas Django dan Flask pada sebagian besar skenario, walaupun hasil semacam ini bergantung pada rancangan pengujiannya. Vue.js tetap dipakai karena model pemrogramannya yang deklaratif dan reaktif \citep*{vuejs2026introduction} cocok untuk kalkulator yang menghitung ulang takaran setiap kali pilihan pesanan berubah. Perubahan ini diputuskan pada 30 September 2026; konsekuensinya, aplikasi memerlukan koneksi ke server, dan mode luring menjadi saran pengembangan pada Bab~\ref{bab:simpulan}.

Berdasarkan uraian tersebut, penelitian ini merancang dan membangun Takar (dari kata \emph{takar}, yaitu mengukur dengan takaran), aplikasi web manajemen resep dan standardisasi operasional kedai kopi dengan rekomendasi berbasis kecerdasan buatan. Takar menyimpan resep standar dari arahan beserta parameter seduh dan SOP-nya, menyediakan kalkulator takaran berbasis batasan dan panduan seduh bertahap, memotong stok bahan baku secara otomatis dari produksi yang dicatat, dan menyediakan tiga modul AI, yaitu rekomendasi pasangan menu dengan FP-Growth, rekomendasi \textit{restock} berbasis peramalan permintaan, dan sistem pakar koreksi parameter seduh. Kontribusi penelitian ini tidak terletak pada kebaruan algoritma secara mutlak, tetapi pada integrasi komponen-komponen tersebut dalam satu aplikasi untuk bar kopi dan pada evaluasi yang terkendali: penskalaan dibandingkan dengan penskalaan linier; modul pasangan menu, peramalan, dan \textit{restock} dibandingkan dengan pembanding sederhana; sistem pakar diuji dengan tabel keputusan dan divalidasi pada data konsumen dan resep seduh nyata; dan hasil pada data sampel fiktif dilaporkan berdampingan dengan hasil pada data nyata.

Di atas ketiga modul tersebut, Takar menyediakan asisten tanya-jawab barista berbasis model bahasa lokal. Pencarian dokumen melalui \textit{retrieval-augmented generation} (RAG) memberi konteks resep, SOP, dan aturan seduh yang dapat dirujuk ketika jawaban disusun \citep{lewis2020retrievalaugmented}, sedangkan pemanggilan alat memungkinkan perhitungan diserahkan kepada fungsi aplikasi \citep{schick2023toolformer}. Model bahasa dapat menghasilkan pernyataan yang tampak meyakinkan tanpa dukungan sumber \citep{ji2023survey}; karena itu jawaban numerik diperiksa terhadap sumber dan alat, dan hasilnya dibandingkan dengan model tanpa konteks pada tolok ukur tanya-jawab yang sama. Tolok ukur yang disiapkan dengan bantuan AI ini belum ditinjau manual oleh peneliti maupun barista praktik, sehingga hasilnya merupakan ukuran teknis internal.

\section{Rumusan Masalah}
\label{sec:rumusan-masalah}

Rumusan masalah penelitian ini dikembangkan dari tiga rumusan masalah pada arahan penelitian. Rumusan tentang sistem manajemen resep digital menjadi RM1, dan rumusan tentang kalkulator resep dinamis menjadi RM2. Rumusan tentang arsitektur PWA tidak lagi dibahas karena perubahan teknologi yang dijelaskan pada Subbab~\ref{sec:latar-belakang}, tetapi tuntutan akses yang cepat dan responsif di area bar tetap diperiksa melalui pengujian kinerja pada RM4. RM3 ditambahkan untuk modul rekomendasi berbasis AI. Dengan demikian, penelitian ini menjawab empat pertanyaan berikut.

\begin{enumerate}[label=\textbf{RM\arabic*}, leftmargin=*, labelsep=1em]
\item Bagaimana merancang dan membangun aplikasi web manajemen resep dan standardisasi operasional kedai kopi yang mendokumentasikan takaran (rasio seduh, gramasi, dan waktu ekstraksi), SOP, serta inventaris bahan baku secara terintegrasi?
\item Bagaimana merancang kalkulator takaran dinamis yang menghitung ulang kebutuhan bahan baku untuk variasi pesanan dengan tetap memenuhi batasan praktis (\textit{shot} utuh, kapasitas cangkir, dan rasio seduh), dan bagaimana hasilnya dibandingkan dengan penskalaan linier?
\item Bagaimana menerapkan kecerdasan buatan untuk rekomendasi pasangan menu, rekomendasi pengadaan (\textit{restock}) bahan baku berbasis peramalan permintaan, koreksi parameter seduh, dan asisten tanya-jawab barista berbasis model bahasa lokal dengan RAG serta pemanggilan alat, dan seberapa baik kinerjanya pada data sampel fiktif, data nyata, dan tolok ukur tanya-jawab?
\item Bagaimana kelayakan aplikasi ditinjau dari pengujian fungsional \textit{black-box}, kinerja, dan penerimaan pengguna?
\end{enumerate}

\section{Hipotesis}
\label{sec:hipotesis}

RM2 dan RM3 dijawab melalui pengujian lima hipotesis berikut, sedangkan RM1 dan RM4 dijawab secara deskriptif melalui rancangan, implementasi, dan hasil pengujian aplikasi.

\begin{enumerate}[label=\textbf{H\arabic*}, leftmargin=*, labelsep=1em]
\item (Penskalaan) Algoritma penskalaan berbasis batasan tidak menghasilkan pelanggaran batasan praktis, yaitu \textit{shot} pecahan, luapan cangkir lebih dari 5\%, dan penyimpangan rasio seduh lebih dari 0,2, pada seluruh skenario uji, sedangkan penskalaan linier menghasilkan pelanggaran.
\item (Pasangan menu) \textit{Hit rate} pada lima rekomendasi teratas (HR@5) FP-Growth lebih tinggi daripada HR@5 \textit{baseline} popularitas pada data uji IBM dan Bread Basket, dengan interval kepercayaan \textit{bootstrap} 95\% untuk selisih berpasangannya berada di atas nol.
\item (Peramalan) Model peramalan terpilih menghasilkan \textit{weighted absolute percentage error} (WAPE) yang lebih rendah daripada \textit{seasonal naive} pada evaluasi \textit{rolling origin} dengan data Maven dan ihelon.
\item (\textit{Restock}) Kebijakan \textit{restock} berbasis peramalan menghasilkan \textit{fill rate} yang lebih tinggi daripada ROP statis pada simulasi persediaan.
\item (Asisten model bahasa lokal) Pada tolok ukur tanya-jawab yang sama, kondisi model bahasa dengan RAG dan pemanggilan alat menghasilkan akurasi jawaban lebih tinggi serta tingkat halusinasi angka lebih rendah daripada model tanpa konteks (\textit{closed-book}).
\end{enumerate}

Luapan cangkir berarti volume cairan melebihi 1,05 kali kapasitas cangkir, sedangkan penyimpangan rasio berarti selisih mutlak antara rasio seduh aktual dan rasio target setelah pembulatan lebih dari 0,2. H1 didukung bila algoritma berbasis batasan tidak menghasilkan satu pun pelanggaran pada seluruh kombinasi resep dan variasi pesanan yang diuji, sedangkan penskalaan linier menghasilkan sedikitnya satu pelanggaran. H2 dan H3 dinilai terpisah untuk setiap \textit{dataset}. Pada H2, selisih HR@5 dihitung berpasangan pada keranjang uji yang sama, dan hipotesis didukung bila batas bawah interval kepercayaannya lebih besar dari nol. Pada H3, model terpilih adalah model dengan WAPE keseluruhan terendah pada data pengembangan Maven, sehingga perbandingan pada data ihelon menjadi pemeriksaan yang tidak dipengaruhi proses pemilihan model. H4 didukung bila proporsi permintaan yang terpenuhi dari stok (\textit{fill rate}) pada kebijakan berbasis peramalan lebih tinggi daripada pada kebijakan ROP statis dalam simulasi yang sama. H5 dinilai pada model utama Gemma 4 E4B dan didukung hanya bila dua syaratnya terpenuhi sekaligus; Gemma 4 E2B dilaporkan sebagai pembanding tambahan. Hasil yang tidak mendukung hipotesis tetap dilaporkan apa adanya pada Bab~\ref{bab:hasil}.

\section{Ruang Lingkup}
\label{sec:ruang-lingkup}

Ruang lingkup penelitian ini dibatasi sebagai berikut.

\begin{enumerate}[label=\arabic*.]
\item \textbf{Fitur.} Aplikasi Takar mencakup (a) manajemen resep standar per kategori beserta komposisi per ukuran, parameter mutu, dan versi resep; (b) kalkulator takaran dinamis untuk variasi ukuran, jumlah cangkir, tingkat kemanisan (0, 50, 100, dan 150\%), tambahan \textit{shot} (0--3), penyajian panas atau dingin, penggantian susu atau biji kopi, dan dosis V60, lengkap dengan peringatan batasan dan perbandingan dengan penskalaan linier; (c) panduan seduh bertahap dengan pengatur waktu dan target tuang yang ikut diskalakan; (d) SOP per resep dan SOP umum beserta daftar periksa pelaksanaannya; (e) inventaris bahan baku berupa buku besar pergerakan stok yang mencatat penerimaan, \textit{stock opname}, bahan terbuang, dan pemakaian otomatis dari produksi; (f) modul rekomendasi AI untuk pasangan menu, peramalan permintaan dan rekomendasi \textit{restock}, serta sistem pakar koreksi parameter seduh; (g) asisten tanya-jawab lokal dengan pencarian dokumen dan pemanggilan layanan aplikasi; serta (h) dasbor dan manajemen pengguna.
\item \textbf{Basis pengetahuan.} Tujuh resep dari arahan disalin apa adanya dalam tiga kategori arahan (\textit{Espresso-Based}, \textit{Manual Brew}, dan \textit{Non-Coffee \& Signature Drinks}), yaitu Caffè Latte, Cappuccino, Ice Brown Sugar Shaken Espresso, Americano/Long Black, V60 dengan biji \textit{single origin} Gayo, Toraja, Flores, dan Ethiopia, Iced Matchagato/Matcha Latte, serta Signature Pandan Latte. Resep tambahan cokelat panas, chai latte, dan teh seduh ditempatkan pada kategori keempat buatan penulis, yaitu Lainnya, sedangkan espreso \textit{double shot} masuk kategori \textit{Espresso-Based}. Resep tambahan tersebut serta setiap angka yang tidak tercantum dalam arahan, seperti ukuran gelas, dosis espreso, massa es, stok awal, dan waktu tunggu pemasok, diperlakukan sebagai asumsi yang didokumentasikan.
\item \textbf{Pengguna.} Aplikasi memiliki tiga peran dengan kontrol akses berbasis peran, yaitu manajer, \textit{head barista}, dan barista. Barista dapat melihat resep, SOP, kalkulator, dan rekomendasi, mencatat produksi, serta menjalankan daftar periksa SOP. \textit{Head barista} dan manajer juga dapat mengubah resep, SOP, dan data bahan serta mencatat penerimaan, \textit{stock opname}, dan bahan terbuang, sedangkan hanya manajer yang dapat mengelola pengguna, melatih ulang model AI, dan menonaktifkan resep.
\item \textbf{Teknologi.} Sisi \textit{backend} dibangun dengan Python, FastAPI, SQLAlchemy \citep*{bayer2026sqlalchemy}, dan SQLite sebagai REST API; sisi \textit{frontend} dibangun dengan Vue.js sebagai \textit{single-page application} yang responsif untuk tablet di area bar dan ponsel. Autentikasi memakai JSON Web Token (JWT), kata sandi disimpan dengan \textit{hash} Argon2, dan modul AI memakai pandas, scikit-learn, statsmodels, serta mlxtend. Asisten menggunakan pencarian BM25 dan model Gemma 4 yang dijalankan secara lokal melalui Ollama; ketika model tidak tersedia, aplikasi menampilkan hasil pencarian dengan penanda mode. Versi setiap perangkat lunak dicatat pada Bab~\ref{bab:hasil}.
% Periode data: IBM dan Maven dari app/SPEC.md bagian 7.1; Bread Basket, ihelon, dan UC Davis dari
% dataset/raw/*/DATASHEET.md (dicocokkan ulang dengan metadata Kaggle/DataCite pada 30 September 2026).
% Kunci dataset: sarwar2018transactions (berkas yang dipakai E1) dan mittal2020bread (unggahan ulang
% berlabel CC0 yang menyebut toko roti di Edinburgh), isaienkov2025coffee (versi 21),
% ristenpart2023consumer (Dryad, versi 4), dan tatman2017nutrition (versi 2).
\item \textbf{Data.} Data pengembangan adalah IBM Coffee Shop Sample Data \citep*{chang2019coffee} berupa struk tiga gerai pada April 2019 untuk rekomendasi pasangan menu, serta Maven Coffee Shop Sales \citep*{mavenanalytics2023coffee} berupa transaksi tiga gerai pada Januari--Juni 2023 untuk peramalan dan simulasi \textit{restock}. Kedua \textit{dataset} tersebut adalah data sampel fiktif: IBM menyebut datanya berasal dari jaringan kedai kopi fiktif, Maven Analytics menyebut datanya catatan transaksi kedai kopi fiktif, dan audit \textit{dataset} dalam penelitian ini menemukan bahwa data Maven memakai katalog produk IBM dan pola bulanan yang berulang. Validasi eksternal memakai data nyata, yaitu transaksi toko roti-kafe The Bread Basket di Edinburgh pada Oktober 2016--April 2017 \citep*{sarwar2018transactions, mittal2020bread} untuk pasangan menu, dengan catatan lisensinya tidak jelas; data penjualan mesin penjual kopi otomatis Coffee Sales dari pengguna Kaggle ihelon (lisensi CC0) pada Maret 2024--Maret 2025 \citep*{isaienkov2025coffee} untuk peramalan dan simulasi \textit{restock}; data preferensi konsumen kopi hitam berlisensi CC0 dari University of California, Davis \citep*{ristenpart2023consumer} yang menyertai studi \citet*{cotter2020consumer} serta tabel resep seduh tetes \citet*{batali2020brew} untuk sistem pakar; dan \textit{dataset} menu Starbucks \citep*{tatman2017nutrition} untuk pola \textit{shot} per ukuran. Produk pada \textit{dataset} dipetakan ke resep aplikasi agar ramalan dapat diturunkan menjadi kebutuhan bahan (Subbab~\ref{sec:dataset}).
\item \textbf{Pengujian.} Pengujian meliputi pengujian fungsional \textit{black-box}, pengukuran kinerja \textit{endpoint} API pada komputer pengembangan, dan tolok ukur internal asisten yang membandingkan pencarian saja, model tanpa konteks, model dengan RAG, serta model dengan RAG dan alat. Uji penerimaan pengguna (\textit{user acceptance testing}, UAT) bersama barista kedai mitra direncanakan, tetapi kedai dan peserta belum ditetapkan; skenario tugas dan kuesioner \textit{System Usability Scale} (SUS) telah disiapkan, sedangkan hasilnya baru dapat dilaporkan setelah UAT dilaksanakan.
\item \textbf{Hal di luar cakupan.} Aplikasi tidak menangani transaksi penjualan dan pembayaran (\textit{point of sale}), tidak menyediakan kemampuan PWA maupun mode luring, dan tidak menyinkronkan data antargerai. Aplikasi tidak terhubung dengan timbangan, refraktometer, atau mesin espreso, sehingga parameter seduh dimasukkan secara manual. Model AI tidak dilatih dengan data penjualan kedai mitra, dan data stok pada aplikasi demonstrasi adalah data contoh. Penerapan (\textit{deployment}) dibatasi pada satu komputer sebagai server lokal (\textit{localhost}).
\end{enumerate}

\section{Tujuan dan Manfaat}
\label{sec:tujuan-manfaat}

\subsection{Tujuan}
\label{subsec:tujuan}

Tujuan penelitian ini, masing-masing menjawab satu rumusan masalah, adalah sebagai berikut.

\begin{enumerate}[label=\arabic*.]
\item Merancang dan membangun aplikasi web Takar yang mendokumentasikan takaran (rasio seduh, gramasi, dan waktu ekstraksi), SOP, dan inventaris bahan baku kedai kopi secara terintegrasi (RM1).
\item Merancang algoritma kalkulator takaran dinamis berbasis batasan dan mengukur pelanggaran batasan praktisnya dibandingkan dengan penskalaan linier, termasuk menguji H1 (RM2).
\item Menerapkan rekomendasi pasangan menu, rekomendasi \textit{restock} berbasis peramalan permintaan, sistem pakar koreksi parameter seduh, dan asisten model bahasa lokal, lalu mengukur kinerjanya pada data sampel fiktif, data nyata, serta tolok ukur tanya-jawab, termasuk menguji H2--H5 (RM3).
\item Menilai kelayakan aplikasi melalui pengujian fungsional \textit{black-box}, pengukuran kinerja, dan uji penerimaan pengguna (RM4).
\end{enumerate}

\subsection{Manfaat}
\label{subsec:manfaat}

Manfaat yang diharapkan dari penelitian ini adalah sebagai berikut.

\begin{enumerate}[label=\arabic*.]
\item \textbf{Bagi kedai kopi dan barista.} Resep standar, parameter seduh, dan SOP tersedia di satu tempat yang sama untuk setiap \textit{shift}, sehingga takaran tidak bergantung pada ingatan barista tertentu. Kalkulator takaran menggantikan hitung ulang manual untuk variasi pesanan dan memberi peringatan bila pesanan melanggar batasan praktis, sedangkan panduan seduh bertahap membantu staf baru mengikuti prosedur yang sama dengan barista berpengalaman. Asisten lokal memudahkan pencarian prosedur beserta sumbernya, dengan hasil perhitungan yang berasal dari layanan aplikasi.
\item \textbf{Bagi manajer kedai.} Stok bahan baku berkurang otomatis dari produksi yang dicatat, sehingga posisi stok dan riwayat pergerakannya dapat dipantau tanpa rekap terpisah. Manajer memperoleh rekomendasi \textit{restock} yang disertai angka dasar perhitungannya, prakiraan permintaan, dan pasangan menu yang dapat dijadikan bahan penawaran paket, dengan hak akses yang dibatasi sesuai peran.
\item \textbf{Bagi akademik.} Penelitian ini menyumbangkan algoritma penskalaan resep berbasis batasan untuk bar kopi beserta evaluasinya terhadap penskalaan linier, serta protokol evaluasi rekomendasi AI yang memakai pembagian data berbasis waktu, pembanding sederhana, dan interval kepercayaan \textit{bootstrap}. Perbandingan empat kondisi asisten lokal memberi contoh pengukuran berbasis sumber, alat, dan kondisi tanpa konteks, dengan batas bahwa tolok ukurnya belum divalidasi manual. Skrip evaluasi mencatat \textit{seed} acak, jumlah baris data, parameter, dan waktu pembuatan hasil sehingga hasilnya dapat direproduksi, dan pemisahan yang tegas antara data sampel fiktif dan data nyata dapat menjadi contoh pelaporan yang jujur bagi penelitian sejenis.
\end{enumerate}

\section{Metode Penelitian}
\label{sec:metode-ringkas}

Penelitian ini merupakan penelitian rancang bangun perangkat lunak yang dilengkapi evaluasi eksperimental algoritma. Pengembangan aplikasi mengikuti siklus hidup pengembangan perangkat lunak dengan model Agile/Scrum sesuai arahan. Scrum adalah kerangka kerja ringan yang menghasilkan nilai melalui solusi adaptif: pekerjaan diurutkan dalam \textit{Product Backlog}, sebagian dikerjakan menjadi \textit{Increment} dalam \textit{Sprint} berdurasi paling lama satu bulan, lalu hasilnya diperiksa dan disesuaikan \citep*{schwaber2020scrum}. Karena dikerjakan oleh satu mahasiswa, Scrum diterapkan dalam bentuk yang disesuaikan; modifikasi Scrum di luar konteks tradisional juga lazim dilaporkan \citep*{hron2022why}. Keterkaitan setiap rumusan masalah dengan metode, instrumen, dan keluarannya diringkas pada Tabel~\ref{tab:rm-metode}, sedangkan tahapan penelitiannya diringkas sebagai berikut dan dibahas rinci pada Bab~\ref{bab:metode}.

\begin{table}[!htb]
\centering
\caption{Keterkaitan Rumusan Masalah, Metode, Instrumen, dan Keluaran Penelitian}
\label{tab:rm-metode}
\small
\begin{tabular}{@{}>{\raggedright\arraybackslash}p{1.0cm}>{\raggedright\arraybackslash}p{4.3cm}>{\raggedright\arraybackslash}p{4.0cm}>{\raggedright\arraybackslash}p{3.9cm}@{}}
\toprule
RM & Metode & Instrumen & Keluaran \\
\midrule
RM1 & Agile/Scrum: analisis kebutuhan, perancangan (UML, ERD, REST API), dan implementasi per \textit{Sprint} & Arahan penelitian, spesifikasi sistem, diagram UML dan ERD, pengujian otomatis & Aplikasi web Takar beserta basis pengetahuan resep dan SOP \\
\addlinespace
RM2 & Perancangan algoritma penskalaan berbasis batasan; eksperimen faktorial penuh terhadap penskalaan linier; pemeriksaan pola \textit{shot} pada \textit{dataset} menu Starbucks & Skrip evaluasi penskalaan dan basis pengetahuan resep & Jumlah dan laju pelanggaran batasan, waktu komputasi, keputusan H1 \\
\addlinespace
RM3 & FP-Growth, peramalan, simulasi \textit{restock}, validasi sistem pakar, serta perbandingan empat kondisi asisten lokal & \textit{Dataset} IBM, Maven, Bread Basket, ihelon, UC Davis, Batali, dan tolok ukur tanya-jawab internal & HR@$K$, MRR@5, WAPE, MASE, \textit{fill rate}, akurasi jawaban, angka tanpa dasar; keputusan H2--H5 \\
\addlinespace
RM4 & Pengujian \textit{black-box}, pengukuran latensi API, dan UAT & Kasus uji otomatis, skrip pengukuran latensi, skenario tugas, dan kuesioner SUS & Laporan pengujian, latensi p50/p95, dan skor SUS setelah UAT dilaksanakan \\
\bottomrule
\end{tabular}
\end{table}

\begin{enumerate}[label=\arabic*.]
\item \textbf{Analisis kebutuhan.} Alur kerja nyata di bar, seperti ekstraksi espreso, rasio V60, dan manajemen stok, diterjemahkan dari arahan penelitian menjadi spesifikasi sistem, \textit{user story}, dan kebutuhan nonfungsional.
\item \textbf{Perancangan sistem.} Tahap ini menghasilkan arsitektur klien--server, diagram UML, rancangan basis data relasional dan ERD, rancangan \textit{endpoint} REST API, dan antarmuka yang ramah layar sentuh, serta rancangan algoritma penskalaan, \textit{restock}, aturan sistem pakar, dan asisten lokal dengan pencarian dokumen serta alat.
\item \textbf{Implementasi bertahap.} \textit{Backend}, \textit{frontend}, dan modul AI dibangun per \textit{Sprint}, dan pengujian otomatis ditulis bersamaan dengan kodenya.
\item \textbf{Evaluasi algoritma pada \textit{dataset}.} Penskalaan dievaluasi pada seluruh kombinasi faktorial resep dan variasi pesanan terhadap penskalaan linier. Pasangan menu dievaluasi dengan pembagian data berbasis waktu dan skema \textit{leave-one-out} per keranjang memakai HR@$K$, \textit{precision}@5, dan \textit{mean reciprocal rank} (MRR@5) terhadap pembanding popularitas dan kemunculan bersama (\textit{co-occurrence}). Peramalan dievaluasi dengan \textit{mean absolute error} (MAE), \textit{root mean squared error} (RMSE), WAPE, dan \textit{mean absolute scaled error} (MASE) pada skema \textit{rolling origin} berhorizon tujuh hari, yaitu titik awal ramalan digeser beberapa kali dan model dilatih ulang pada setiap titik \citep*{hewamalage2022forecast}. Kebijakan \textit{restock} dibandingkan melalui simulasi harian, dan sistem pakar diuji dengan tabel keputusan serta divalidasi pada data konsumen dan resep seduh nyata. Asisten diuji pada pertanyaan yang sama dalam empat kondisi, termasuk model tanpa konteks sebagai pembanding H5. Rancangan eksperimen dibahas pada Subbab~\ref{sec:rancangan-eksperimen}.
\item \textbf{Pengujian aplikasi.} Pengujian fungsional \textit{black-box} memakai teknik \textit{equivalence partitioning} (EP) dan \textit{boundary value analysis} (BVA) \citep*{myers2012art} yang dijalankan sebagai pengujian API otomatis dan pengujian antarmuka \textit{end-to-end}. Kinerja diukur sebagai latensi persentil ke-50 dan ke-95 (p50/p95) pada \textit{endpoint} utama. UAT, yaitu validasi perangkat lunak dalam situasi nyata oleh calon penggunanya \citep*{otaduy2017user}, direncanakan memakai skenario tugas dan kuesioner SUS sepuluh butir \citep*{brooke1996sus}. Rancangan pengujian dibahas pada Subbab~\ref{sec:rancangan-pengujian}.
\end{enumerate}

\section{Sistematika Penulisan}
\label{sec:sistematika}

Skripsi ini terdiri atas lima bab dengan rincian sebagai berikut.

\textbf{BAB 1 PENDAHULUAN}

Bab ini membahas latar belakang penelitian, rumusan masalah, hipotesis, ruang lingkup, tujuan dan manfaat, metode penelitian, serta sistematika penulisan.

\textbf{BAB 2 TINJAUAN REFERENSI}

Bab ini membahas teori yang mendasari penelitian, yaitu operasional kedai kopi dan resep standar, ilmu penyeduhan kopi, penskalaan resep, pengendalian persediaan, sistem berbasis web beserta keamanannya, kecerdasan buatan dan sistem rekomendasi, aturan asosiasi, peramalan permintaan, sistem pakar, model bahasa lokal dan RAG, evaluasi sistem rekomendasi, Scrum, UML, dan pengujian perangkat lunak. Bab ini juga memuat penelitian terkait dan kerangka pemikiran.

\textbf{BAB 3 METODE PENELITIAN}

Bab ini membahas desain dan tahapan penelitian, analisis kebutuhan, perancangan sistem (arsitektur, diagram UML, ERD, peta situs, dan \textit{endpoint} API), perancangan algoritma penskalaan, \textit{restock}, aturan sistem pakar dan asisten lokal, \textit{dataset} beserta pemetaan produknya, rancangan eksperimen dan pengujian, serta jadwal penelitian.

\textbf{BAB 4 HASIL DAN PEMBAHASAN}

Bab ini membahas lingkungan dan hasil implementasi, hasil pengujian fungsional, hasil evaluasi penskalaan, pasangan menu, peramalan, \textit{restock}, dan sistem pakar pada data sampel fiktif dan data nyata, tolok ukur internal asisten lokal, hasil pengujian kinerja dan status UAT, serta pembahasan hasil terhadap hipotesis dan keterbatasannya.

\textbf{BAB 5 SIMPULAN DAN SARAN}

Bab ini membahas simpulan yang menjawab setiap rumusan masalah dan saran untuk pengembangan penelitian selanjutnya.
